\chapter{分歧怎样改变三国史的读法}
\label{app:debates}

分歧不只是脚注里的障碍。它会改变我们如何理解选择、责任和国家能力。本附录
不裁决所有异说，而是列出几个会实质改变叙事尺度的地方。

\section{禅代：礼仪记录与实际选择}

\sourcebook{后汉书}、\sourcebook{三国志}和\sourcebook{晋书}都保留了
汉魏、魏晋转换的制度语言。它们可以支持\eventref{EH-0220-HAN-ABDICATION}
{汉魏禅代}和\eventref{TK-0265-WEI-ABDICATION}{魏晋禅代}的次序，却不让
读者直接进入诏书起草、禁军部署和群臣被迫或自愿的程度。将“禅让”只读作
和平礼仪，会隐去事先取得军政控制的条件；只读作暴力夺权，又会忽略礼仪、
官僚与文本如何使转换可以被日常行政承接。

\section{夷陵、北伐与数字的诱惑}

\eventref{TK-0222-YILING}{夷陵}、\eventref{TK-0228-JIETING}{街亭}和
\eventref{TK-0234-WUZHANG}{五丈原}在后世叙事中常有清晰的英雄、谋略与
伤亡数字。正史和裴注足以确定若干行动与结果，未必足以确证营垒长度、火攻
步骤、军队规模或单一责任人。把这些细节降为不确定项，不会削弱解释，反而
使山道、粮运、饮水、守险和盟友关系能重新进入战争的原因链。

\section{边疆不等于空白}

辽东、高句丽、倭、南中和羌胡各部在中央史书中常以一场征伐或一次册封出现。
\eventref{TK-0238-SIMA-YI-LIAODONG}{平辽东}、毌丘俭破丸都和魏受倭使，
能说明魏军与魏廷抵达过哪些节点；它们不能自动证明持续直辖，亦不能使族名
等于固定的现代身份。区域史、山城考古、地名研究和文书的作用，是让中央
叙事的范围受到空间和地方经验的检验。

\section{吴简与石经的局部亮度}

走马楼吴简不应被缩小为“没有用的碎片”，也不应被放大为完整的吴国统计局；
它最有力地改变了我们看待家庭、田租和吏役的方式。正始石经同样如此：
\eventref{TK-0240-ZHENGSHI-STONE-CLASSICS}{立石经}能说明洛阳的官学和
书写权威，不能证明一统的学术或政治认同。材料的地域与保存偶然性，是结论
尺度的一部分。

\section{统一的速度}

\eventref{TK-0279-JIN-INVADES-WU}{晋军伐吴}和
\eventref{TK-0280-SUN-HAO-SURRENDER}{孙皓出降}清楚结束了三国最高政权的
并立。把它描述成江南立刻与北方同质化，却会越过我们拥有的证据。西晋接收
的是孙吴州郡、簿籍、官民与边防网络；泰始律和新的任命为接收提供制度语言，
地方执行、家庭负担和区域市场的改变仍须在后续时代的材料中逐项追踪。
